\scanpage{40}
% archive sheet number: 39
Donc \deleted{\unsure{la|les} \illegible{1}} on trouve

\uline{Proposition} $A$, \inserted{$h$,} $W,e,f$, \inserted{$h=ef$,} comme plus haut, $G_0$ un groupe de B.T: \inserted{de hauteur $h$} sur \deleted{$A$} $k$ sur lequel $A$ opère, \deleted{Alors \unsure{$F$|si}} $M$ son module de Dieudonné (module libre de rang $h$ sur $W$). Alors $F_{G_0}^f$ est de la forme $\varpi_{G_0}$, où $\varpi\in A$, et on a $\dim G_0=\operatorname{long} A/\varpi A=v(\varpi)$.

Alors \deleted{conditions équivalentes}

Conditions équivalentes

(i) $\dim G_0=1$

(ii) $F_{G_0}^f=\pi_{G_0}$, $\pi$ une \uline{uniformisante} de $A$ (uniqt déterminée !)

(iii) Le module de Dieudonné $M$ de $G_0$ \inserted{(sur $W$)}, avec $F_M$ et action de $A$, est donné par un triple $(\Omega,\unsure{\nu_0|\nu},\pi)$, où $\Omega$ est un module inv. sur $A$, $0\leq\nu\leq f-1$, $\pi$ une uniformisante de $A$ par
\[
\left\{\begin{aligned}
M&=\Omega\otimes_{\mathbb Z_p}W=\Omega\otimes_A(A\otimes_{\mathbb Z_p}W)\qquad(\text{isom de }A\otimes_{\mathbb Z_p}W\text{-mod})\\
F_M&=\mathrm{id}_\Omega\otimes_{\mathbb Z_p}\sigma+\operatorname{pr}_{\nu}[(\pi-1)\otimes\sigma]\deleted{\unsure{\operatorname{pr}_{\nu}}}
\end{aligned}\right.
\]
où $\operatorname{pr}_\nu$ est la \deleted{facteur} projection de $M$ sur son $\nu$-ième facteur \inserted{$M_\nu$} (correspondant à la décomposition $A\otimes_{\mathbb Z_p}W\simeq A^{[0,f-1]}$)

% slanted
\margin{\uline{NB} Tous ces groupes (pour $\nu$ \unsure{variable}) sont \unsure{déduits|isomorphes}, de l'un d'entre eux (disons pour $\nu=0$) par \unsure{transformation|transport|transformations} à l'aide des $f$ automorphismes $\sigma^i$ de $k=\mathbb F_q=\mathbb F_{p^f}$.}

\uline{NB} Revenons au cas général où on ne suppose pas néc. $\dim G_0=1$ \inserted{(mais $G_0$ de hauteur $h$ $\ldots$)}. Soit $G'_0$ défini comme $G_0$ par $\Omega'$, \deleted{\illegible{1}} les $r'_\nu$ $(0\leq\nu\leq f-1)$ $u'\in A^*$. On cherche les hom compatibles avec l'action de $A$ de $G_0$ dans $G'_0$, i.e. les hom de $A\otimes_{\mathbb Z_p}W$-modules de $M'$ en $M$, compatibles avec $F_M,F_{M'}$. On trouve \unsure{facilt|facilemt|facilement} que pour qu'il existe un tel hom qui \deleted{\unsure{soit}} $\ne 0$ ($\Longleftrightarrow$ une isogénie !), il f. et suffit qu'on ait $\varpi=\varpi'$ i.e. $F_{G_0}^f=F_{G'_0}^f$ (dans $A$). Dans ce cas, \deleted{\unsure{les}} le gpe des hom en question, comme $A$-module, est isom au \deleted{\unsure{$E'$}} sous-module $H\subset\operatorname{Hom}(\Omega',\Omega)$ formé des $v_0:\Omega'\to\Omega$ tels qu'on ait
\[
v_0\pi^{(\sum_0^\nu r_i)-(\sum_0^\nu r'_i)}\in\operatorname{Hom}(\Omega',\Omega)\quad\inserted{\text{pour tt }0\leq\nu\leq f-2},\quad\text{i.e.}\quad \operatorname{val}(v_0)+\sum_0^\nu r_i-\sum_0^\nu r'_i\geq 0
\]
